\documentclass[10pt,a4paper]{amsart}
\usepackage{amsmath,amssymb,amsthm,graphicx}
% Retained minimal preamble for the Stacks Project chapter flat.tex (More on Flatness).
% Verbatim excerpt of the chapter preamble: Xy-pic setup (native preamble lines 14-18) and the native
% theorem declarations with the mathematical macros (native preamble lines 154-217). Copied mechanically;
% nothing added, removed or redefined. The wrapper class and the amsmath/amssymb/amsthm/graphicx base
% packages are supplied by the assembler, so the class/externaldocument part of the native preamble is omitted.
% For commutative diagrams we use Xy-pic
\usepackage[all]{xy}

% We use 2cell for 2-commutative diagrams.
\xyoption{2cell}
% Theorem environments.
%
\theoremstyle{plain}
\newtheorem{theorem}[subsection]{Theorem}
\newtheorem{proposition}[subsection]{Proposition}
\newtheorem{lemma}[subsection]{Lemma}

\theoremstyle{definition}
\newtheorem{definition}[subsection]{Definition}
\newtheorem{example}[subsection]{Example}
\newtheorem{exercise}[subsection]{Exercise}
\newtheorem{situation}[subsection]{Situation}

\theoremstyle{remark}
\newtheorem{remark}[subsection]{Remark}
\newtheorem{remarks}[subsection]{Remarks}

\numberwithin{equation}{subsection}

% Macros
%
\def\lim{\mathop{\mathrm{lim}}\nolimits}
\def\colim{\mathop{\mathrm{colim}}\nolimits}
\def\Spec{\mathop{\mathrm{Spec}}}
\def\Hom{\mathop{\mathrm{Hom}}\nolimits}
\def\Ext{\mathop{\mathrm{Ext}}\nolimits}
\def\SheafHom{\mathop{\mathcal{H}\!\mathit{om}}\nolimits}
\def\SheafExt{\mathop{\mathcal{E}\!\mathit{xt}}\nolimits}
\def\Sch{\mathit{Sch}}
\def\Mor{\mathop{\mathrm{Mor}}\nolimits}
\def\Ob{\mathop{\mathrm{Ob}}\nolimits}
\def\Sh{\mathop{\mathit{Sh}}\nolimits}
\def\NL{\mathop{N\!L}\nolimits}
\def\CH{\mathop{\mathrm{CH}}\nolimits}
\def\proetale{{pro\text{-}\acute{e}tale}}
\def\etale{{\acute{e}tale}}
\def\QCoh{\mathit{QCoh}}
\def\Ker{\mathop{\mathrm{Ker}}}
\def\Im{\mathop{\mathrm{Im}}}
\def\Coker{\mathop{\mathrm{Coker}}}
\def\Coim{\mathop{\mathrm{Coim}}}

% Boxtimes
%
\DeclareMathSymbol{\boxtimes}{\mathbin}{AMSa}{"02}

%
% Macros for moduli stacks/spaces
%
\def\QCohstack{\mathcal{QC}\!\mathit{oh}}
\def\Cohstack{\mathcal{C}\!\mathit{oh}}
\def\Spacesstack{\mathcal{S}\!\mathit{paces}}
\def\Quotfunctor{\mathrm{Quot}}
\def\Hilbfunctor{\mathrm{Hilb}}
\def\Curvesstack{\mathcal{C}\!\mathit{urves}}
\def\Polarizedstack{\mathcal{P}\!\mathit{olarized}}
\def\Complexesstack{\mathcal{C}\!\mathit{omplexes}}
% \Pic is the operator that assigns to X its picard group, usage \Pic(X)
% \Picardstack_{X/B} denotes the Picard stack of X over B
% \Picardfunctor_{X/B} denotes the Picard functor of X over B
\def\Pic{\mathop{\mathrm{Pic}}\nolimits}
\def\Picardstack{\mathcal{P}\!\mathit{ic}}
\def\Picardfunctor{\mathrm{Pic}}
\def\Deformationcategory{\mathcal{D}\!\mathit{ef}}

% --- appended 2026-10-12 (Chief): the project own definitions of the reference/slogan/history
% display environments, copied mechanically from the shared Stacks preamble lines 8-12
% (\usepackage{verbatim} plus the three \newenvironment lines, which the chapter preamble keeps
% verbatim from the project preamble). Nothing added, removed or redefined: these environments
% carry no mathematical content and the project itself typesets them as comments. Without these
% three lines a counted statement whose own text uses them fails with Environment slogan undefined.
\usepackage{verbatim}
\newenvironment{reference}{\comment}{\endcomment}
\newenvironment{slogan}{\comment}{\endcomment}
\newenvironment{history}{\comment}{\endcomment}
\usepackage{cleveref}
\crefname{originallemma}{original lemma}{original lemmas}
\Crefname{originallemma}{Original lemma}{Original lemmas}
\makeatletter
\newlabel{divisors-lemma-base-change-relative-assassin}{{05DC}{source}{}{sourcepointer-1}{}}
\newlabel{divisors-lemma-base-change-relative-assassin@cref}{{[originallemma][0][]05DC}{[1][0][]source}{}{}{}}
\makeatother
\crefname{originallemma}{original lemma}{original lemmas}
\Crefname{originallemma}{Original lemma}{Original lemmas}
\makeatletter
\newlabel{more-morphisms-lemma-etale-makes-quasi-finite-finite-at-point}{{02LK}{source}{}{sourcepointer-2}{}}
\newlabel{more-morphisms-lemma-etale-makes-quasi-finite-finite-at-point@cref}{{[originallemma][0][]02LK}{[1][0][]source}{}{}{}}
\makeatother
\crefname{originallemma}{original lemma}{original lemmas}
\Crefname{originallemma}{Original lemma}{Original lemmas}
\makeatletter
\newlabel{situation-pre-pure}{{05FW}{source}{}{sourcepointer-3}{}}
\newlabel{situation-pre-pure@cref}{{[originallemma][0][]05FW}{[1][0][]source}{}{}{}}
\makeatother
\title{An impurity with a quasi-finite morphism can be replaced by one over an elementary etale neighbourhood}
\author{The Stacks Project Authors}
\date{stacks-project:flat.tex:local-snapshot-20261011; GNU FDL 1.2 or later}
\begin{document}
\maketitle
\noindent\textit{Editorial note.} All delivered mathematical text is the author's own, copied from the pinned Stacks Project chapter file flat.tex; no statement or proof step was rewritten, reordered or removed.\par\medskip
\noindent\textit{Editorial note.} Preamble provenance: the chapter's own retained minimal preamble (sources/stacks-flat/retained-preamble.tex, lines 1-85) is a verbatim excerpt of the native Xy-pic setup and of the native theorem declarations and mathematical macros, followed by the project's own reference/slogan/history display-environment definitions copied mechanically from the shared preamble; the wrapper class and the amsmath/amssymb/amsthm/graphicx base packages are supplied by the assembler, and extra\_packages is left empty because everything needed is already in the retained preamble. No macro was invented or redefined.\par\medskip
\noindent\textit{Editorial note.} Same-chapter pointers are located at their \textbackslash{}label lines inside the pinned chapter file; cross-chapter pointers are located in the project's own tag list staged in the snapshot (sources/stacks-chow/tags.txt), at the verbatim lines carrying their canonical tags.\par\medskip
\noindent\textit{Allowed standard prerequisites (editorial).}\begin{itemize}
\item The situation and the notion of an impurity of a module above a point. Same chapter flat.tex, situation labelled situation-pre-pure (statement, proof line 4646). Fixes the data of an impurity: a morphism with a specialisation and a point of the relative assassin whose closure misses the fibre over the specialisation.
\item Splitting off a finite part of a quasi-finite morphism over an elementary etale neighbourhood. More on Morphisms, label more-morphisms-lemma-etale-makes-quasi-finite-finite-at-point, invoked at proof line 4659. Produces the elementary etale neighbourhood (U, u) -> (S, s) and the open neighbourhood V of v = (t, u) in the base change which is finite over U with v the unique point over u, which is the construction at the heart of the proof.
\item Behaviour of the relative assassin under base change. Divisors, label divisors-lemma-base-change-relative-assassin, invoked at proof lines 4675 and 4687. Transports a relative associated point of the module along the etale and finite base changes, first to the point zeta over xi and then to its image theta over U.
\end{itemize}\medskip
\ifdefined\hypertarget\hypertarget{sourcepointer-1}{}\fi
\noindent\textit{Source reference.} Original lemma 05DC: sources/stacks-chow/tags.txt, source lines 6950--6950. Navigation pointer for \textbackslash{}ref\{divisors-lemma-base-change-relative-assassin\} used inside the counted span (line 4675); the argument that consumes it is retained in the span.\par\medskip
\ifdefined\hypertarget\hypertarget{sourcepointer-2}{}\fi
\noindent\textit{Source reference.} Original lemma 02LK: sources/stacks-chow/tags.txt, source lines 3411--3411. Navigation pointer for \textbackslash{}ref\{more-morphisms-lemma-etale-makes-quasi-finite-finite-at-point\} used inside the counted span (line 4659); the argument that consumes it is retained in the span.\par\medskip
\ifdefined\hypertarget\hypertarget{sourcepointer-3}{}\fi
\noindent\textit{Source reference.} Original lemma 05FW: sources/stacks-flat/pinned-original-source.tex, source lines 4473--4473. Navigation pointer for \textbackslash{}ref\{situation-pre-pure\} used inside the counted span (line 4646); the argument that consumes it is retained in the span.\par\medskip
\begin{lemma}
\label{lemma-quasi-finite-impurity-elementary}
In Situation \ref{situation-pre-pure}.
If there exists an impurity $(g : T \to S, t' \leadsto t, \xi)$
of $\mathcal{F}$ above $s$ with $g$ quasi-finite at $t$, then there
exists an impurity $(g : T \to S, t' \leadsto t, \xi)$ such that
$(T, t) \to (S, s)$ is an elementary \'etale neighbourhood.
\end{lemma}

\begin{proof}
Let $(g : T \to S, t' \leadsto t, \xi)$ be an impurity of
$\mathcal{F}$ above $s$ such that $g$ is quasi-finite at $t$.
After shrinking $T$ we may assume that $g$ is locally of finite type.
Apply
More on Morphisms,
Lemma \ref{more-morphisms-lemma-etale-makes-quasi-finite-finite-at-point}
to $T \to S$ and $t \mapsto s$. This gives us a diagram
$$
\xymatrix{
T \ar[d] & T \times_S U \ar[l] \ar[d] & V \ar[l] \ar[ld] \\
S & U \ar[l]
}
$$
where $(U, u) \to (S, s)$ is an elementary \'etale neighbourhood
and $V \subset T \times_S U$ is an open neighbourhood of $v = (t, u)$
such that $V \to U$ is finite and such that $v$ is the unique point of $V$
lying over $u$. Since the morphism $V \to T$ is \'etale
hence flat we see that there exists a specialization $v' \leadsto v$ such
that $v' \mapsto t'$. Note that $\kappa(t') \subset \kappa(v')$
is finite separable. Pick any point $\zeta \in X_{v'}$ mapping to
$\xi \in X_{t'}$. By
Divisors, Lemma \ref{divisors-lemma-base-change-relative-assassin}
we see that $\zeta \in \text{Ass}_{X_V/V}(\mathcal{F}_V)$.
Moreover, the closure $\overline{\{\zeta\}}$ does not meet
the fibre $X_v$ as by assumption the closure $\overline{\{\xi\}}$
does not meet $X_t$. In other words $(V \to S, v' \leadsto v, \zeta)$
is an impurity of $\mathcal{F}$ above $S$.

\medskip\noindent
Next, let $u' \in U'$ be the image of $v'$ and let
$\theta \in X_U$ be the image of $\zeta$.
Then $\theta \mapsto u'$ and $u' \leadsto u$.
By
Divisors, Lemma \ref{divisors-lemma-base-change-relative-assassin}
we see that $\theta \in \text{Ass}_{X_U/U}(\mathcal{F})$.
Moreover, as $\pi : X_V \to X_U$ is finite we see that
$\pi\big(\overline{\{\zeta\}}\big) = \overline{\{\pi(\zeta)\}}$. Since
$v$ is the unique point of $V$ lying over $u$ we see that
$X_u \cap \overline{\{\pi(\zeta)\}} = \emptyset$ because
$X_v \cap \overline{\{\zeta\}} = \emptyset$. In this way we conclude that
$(U \to S, u' \leadsto u, \theta)$ is an impurity of
$\mathcal{F}$ above $s$ and we win.
\end{proof}

\end{document}
